\documentclass[12pt, letterpaper]{article}

\usepackage{amsmath, amssymb}
\usepackage[hidelinks]{hyperref}

\setlength{\parindent}{0pt}
\setlength{\parskip}{0.5em}

\begin{document}

\vspace*{-1cm}
\begin{center}
    {\Huge\bfseries Definitions}
\end{center}

\section{Machine Learning}

\subsection{Hypothesis Space}

The set of all functions a model can represent, denoted by
$\mathcal{H}$.

\[
    \mathcal{H} \subseteq
    \{h : \mathcal{X} \to \mathcal{Y}\}
\]

Here, $h$ is a hypothesis mapping inputs $\mathcal{X}$
to outputs $\mathcal{Y}$.

\subsection{Inductive Bias}

The assumptions that guide a learning algorithm toward certain
hypotheses, allowing it to generalize beyond its training data.

\textbf{Restriction bias.}
Limits the hypothesis space $\mathcal{H}$ to a smaller set
of possible functions.

\[
    \mathcal{H}_{\text{restricted}} \subseteq \mathcal{H}
\]

A smaller hypothesis space can require less data to generalize,
provided it contains a sufficiently good hypothesis.

\textbf{Preference bias.}
Favors certain hypotheses over others without necessarily
excluding them from $\mathcal{H}$.

For example, weight decay encourages solutions with smaller weights.

\subsection{L2 Norm} % NOTE: THIS PROBABLY SHOULDNT BE HERE!!

The L2 norm, also known as the Euclidean norm, measures
the magnitude (length) of a vector.

For a vector $w \in \mathbb{R}^n$:

\[
    \|w\|_2 = \sqrt{\sum_{i=1}^{n} w_i^2}
\]

For example, given $w = [3,4]$:

\[
    \|w\|_2 = \sqrt{3^2 + 4^2} = 5
\]

The squared L2 norm removes the square root:

\[
    \|w\|_2^2 = \sum_{i=1}^{n} w_i^2
\]

\subsection{L2 Regularization}

A regularization technique that discourages large model weights,
helping reduce overfitting.

Instead of minimizing only prediction loss, we add a penalty
proportional to the squared L2 norm of the weights.

\[
    L_{\text{total}} =
    L_{\text{prediction}} +
    \frac{\lambda}{2}\|w\|_2^2
\]

Where:
\begin{itemize}
    \item $L_{\text{prediction}}$ is the original loss function.
    \item $w$ is the vector of model weights being regularized.
    \item $\lambda \geq 0$ controls regularization strength.
    \item $\|w\|_2^2$ is the sum of squared weights.
\end{itemize}

Larger values of $\lambda$ penalize large weights more strongly.
When $\lambda = 0$, there is no regularization.

\subsubsection*{Effect on Gradient Descent}

The penalty contributes an additional term to the gradient:

\[
    \frac{\partial L_{\text{total}}}{\partial w}
    =
    \frac{\partial L_{\text{prediction}}}{\partial w}
    + \lambda w
\]

Substituting into the gradient descent update:

\[
    w_{\text{new}} =
    w - \eta
    \left(
        \frac{\partial L_{\text{prediction}}}{\partial w}
        + \lambda w
    \right)
\]

Rearranging:

\[
    \boxed{
        w_{\text{new}} =
        (1-\eta\lambda)w -
        \eta\frac{\partial L_{\text{prediction}}}{\partial w}
    }
\]

Thus, each update shrinks the weights by a factor of
$(1-\eta\lambda)$ before applying the prediction-loss gradient.

This is equivalent to weight decay in standard SGD.

\end{document}
